% -------------------------------------------------------------------------
% ÜBERHOLT (31.08.2026). Maßgeblich ist die eingereichte main.tex. Zahlen,
% Zeilennummern und Angaben zum Overleaf-Stand in dieser Datei stammen aus
% der Zeit VOR der Crosswalk-Korrektur vom 29.08.2026 (35 statt 36 Stadt-
% teile, vollständiger Neulauf) und wurden bewusst NICHT nachgezogen: die
% Datei ist als datierter Entwurf Teil der Arbeitsdokumentation.
% -------------------------------------------------------------------------
% =========================================================================
% CODEAUSSCHNITTE KAPITEL 5 -- DATA PREPARATION
% Stand 22.08.2026. Erzeugt aus dem Repositorium sffd_analyse.
%
% REGEL FUER DIESE DATEI: Jede Zeile steht ZEICHENGLEICH so in der Quelldatei.
% Gekuerzt wird ausschliesslich durch WEGLASSEN ganzer Zeilen, nie durch
% Umformulieren. Damit entsteht keine zweite Fassung des Codes, die
% auseinanderlaufen koennte. Wo ein Kommentar zu lang war, steht er gar nicht
% im Ausschnitt -- seine Begruendung gehoert ohnehin in den Fliesstext.
%
% Wird ein Kommentar im Ausschnitt geaendert, ist er ZUERST in der .py-Datei
% zu aendern und von dort zu uebernehmen. Nicht umgekehrt.
%
% Die Zeilennummern in den Beschriftungen beziehen sich auf den Stand
% 22.08.2026. Sie verschieben sich bei jeder Codeaenderung -- vor der Abgabe
% einmal nachziehen.
%
% BESCHRIFTUNGEN: bewusst einzeilig, nur Gegenstand und Fundstelle. Die
% Begruendung gehoert in den Fliesstext, nicht unter den Ausschnitt --
% sonst steht sie zweimal da.
%
% Voraussetzung im Preamble: \lstdefinestyle{pythonstil}{...}
% (Vorschlag siehe Kapitel5_Codeausschnitte.md, Abschnitt 3)
% =========================================================================


% -------------------------------------------------------------------------
% L1  Publikationsversatz des ACS  ->  Abschnitt Zusammenfuehrung
% -------------------------------------------------------------------------
\begin{lstlisting}[style=pythonstil,
  label={lst:acs-versatz},
  caption={Zeitbewusste Zuordnung des ACS-Jahrgangs
           (\texttt{prep/s1\_daten.py}, Z.~306--322 und 337, gekürzt)}]
def acs_snapshot(jahr: int, acs_years: list[int]) -> int:
    return max([a for a in acs_years if a <= int(jahr) - ACS_PUBLIKATIONS_LAG],
               default=min(acs_years))

sffd["acs_year"] = sffd["year"].apply(lambda y: acs_snapshot(y, jahrgaenge))
\end{lstlisting}


% -------------------------------------------------------------------------
% L2  Quotenbildung, Nenner <= 0  ->  Abschnitt Merkmale
% -------------------------------------------------------------------------
\begin{lstlisting}[style=pythonstil,
  label={lst:quoten},
  caption={Anteilswerte aus Zählvariablen
           (\texttt{prep/s1\_daten.py}, Z.~559--584, gekürzt)}]
QUOTEN = [
    ("poverty_rate",     "poverty_below",         "poverty_universe_total"),
    ("bachelor_rate",    "bachelor_degree_count", "education_universe_total"),
    ("vacancy_rate",     "vacant_housing_units",  "total_housing_units"),
    ("pct_pre1940",      "pre1940_count",         "yrbuilt_count"),
    ("pct_residential",  "residential_count",     "parcel_count"),
]

def berechne_quoten(df: pd.DataFrame) -> pd.DataFrame:
    for name, zaehler, nenner in QUOTEN:
        z = pd.to_numeric(df[zaehler], errors="coerce").astype(float)
        n = pd.to_numeric(df[nenner], errors="coerce").astype(float)
        df[name] = (z / n.where(n > 0, np.nan)).round(4)
    return df
\end{lstlisting}


% -------------------------------------------------------------------------
% L3  Kriminalitaetsindex als Location Quotient  ->  Abschnitt Merkmale
%     Der wichtigste Ausschnitt des Kapitels.
% -------------------------------------------------------------------------
\begin{lstlisting}[style=pythonstil,
  label={lst:krimindex},
  caption={Relativer Kriminalitätsindex je Stadtteil und Monat
           (\texttt{prep/s1\_daten.py}, Z.~465--488, gekürzt)}]
    # Rollierendes Fenster, endend im VORMONAT: shift(1) VOR rolling()
    raster["delikte_fenster"] = (
        raster.groupby("neighborhood")["delikte"]
              .transform(lambda s: s.shift(1).rolling(CRIME_FENSTER_MONATE).sum()))
    raster = raster.dropna(subset=["delikte_fenster"])

    stadt = (raster.groupby(["jahr", "monat"])
                   .agg(stadt_delikte=("delikte_fenster", "sum"),
                        stadt_bev=("total_population", "sum")).reset_index())
    raster = raster.merge(stadt, on=["jahr", "monat"], how="left")

    raster["crime_rate_raw"] = (raster["delikte_fenster"]
                                / raster["total_population"] * 1000).round(3)
    raster["crime_index"] = (raster["crime_rate_raw"]
                             / (raster["stadt_delikte"] / raster["stadt_bev"] * 1000)
                             ).round(4)
\end{lstlisting}


% -------------------------------------------------------------------------
% L4  Lag-Bildung  ->  Abschnitt Merkmale
%     ACHTUNG Dopplung: L3 zeigt shift(1) VOR rolling() bereits. Entweder
%     hier nur lag_1/lag_12 zeigen oder im Fliesstext auf L3 verweisen.
% -------------------------------------------------------------------------
\begin{lstlisting}[style=pythonstil,
  label={lst:lags},
  caption={Bildung der Verlaufsmerkmale
           (\texttt{prep/s2\_datensaetze.py}, Z.~269--283, gekürzt)}]
    # LAGS - Strikt rueckwaertsgerichtet und je Stadtteil gebildet, nie ueber
    # Stadtteilgrenzen hinweg. Beim gleitenden Mittel steht `shift(1)` VOR
    # `rolling(3)`: der Wert fuer Monat t verwendet t-1, t-2, t-3, nie t selbst.
    d = d.sort_values(["stadtteil", "jahr_monat"]).reset_index(drop=True)
    g = d.groupby("stadtteil")[ZIELGROESSE]
    d["lag_1"]          = g.shift(1)
    d["lag_12"]         = g.shift(12)
    d["rolling_mean_3"] = g.transform(lambda s: s.shift(1).rolling(3).mean())

    d = d[d["jahr_monat"] >= START].reset_index(drop=True)
\end{lstlisting}


% -------------------------------------------------------------------------
% L5  SERIALISIERUNG  ->  Abschnitt Analyserahmen
%     Auflage Schroeter 10.08.2026: Serialisierung ausdruecklich zeigen.
% -------------------------------------------------------------------------
\begin{lstlisting}[style=pythonstil,
  label={lst:dtypes},
  caption={Vereinheitlichung auf modelltaugliche NumPy-Typen
           (\texttt{prep/s2\_datensaetze.py}, Z.~147--170, gekürzt)}]
def _setze_datentypen(d: pd.DataFrame, merkmale: list[str]) -> pd.DataFrame:
    # notwendig, weil EINE nullable Int64-Spalte genuegt, damit X.to_numpy() ein
    # object-Array liefert
    d = d.copy()
    for c in merkmale:
        d[c] = pd.to_numeric(d[c], errors="coerce").astype("float64")
    for c in ["jahr", "monat", "jahr_monat", ZIELGROESSE,
              "fold", "ist_holdout", EXPOSURE_ROH] + ANZAHLEN:
        if c in d.columns:
            d[c] = pd.to_numeric(d[c], errors="coerce").astype("int64")
    return d
\end{lstlisting}


% -------------------------------------------------------------------------
% L6  Fold-Zuteilung mit doppelter Stratifizierung  ->  Abschnitt Analyserahmen
% -------------------------------------------------------------------------
\begin{lstlisting}[style=pythonstil,
  label={lst:folds},
  caption={Zuteilung der Stadtteile auf Folds und Hold-out
           (\texttt{prep/s2\_datensaetze.py}, Z.~61--89, gekürzt)}]
def ergaenze_aufteilung(daten: pd.DataFrame, versatz: int = 0,
                        selten: pd.Series | None = None) -> pd.DataFrame:
    bev = daten.groupby("stadtteil")[EXPOSURE_ROH].mean()
    if selten is None:
        selten = pd.Series(0, index=bev.index)
    ordnung = (pd.DataFrame({"selten": selten.reindex(bev.index).fillna(0),
                             "bev": bev})
                 .sort_values(["selten", "bev"], ascending=False).index)
    gruppe = {st: (i + versatz) % (N_FOLDS + 1) for i, st in enumerate(ordnung)}

    d = daten.copy()
    d["fold"] = d["stadtteil"].map(gruppe).astype(int)
    d["ist_holdout"] = (d["fold"] == 0).astype(int)
    return d
\end{lstlisting}


% -------------------------------------------------------------------------
% L7  Stufe-2-Baseline  ->  Abschnitt Analyserahmen
%     Belegt Auflage D konstruktiv: unpenalisiert, kein freier Parameter.
% -------------------------------------------------------------------------
\begin{lstlisting}[style=pythonstil,
  label={lst:poisson},
  caption={Stufe-2-Baseline: Poisson-GLM mit Offset
           (\texttt{vorpruefung/v1\_baselines.py}, Z.~74--104, gekürzt)}]
def poisson_glm(train: pd.DataFrame, test: pd.DataFrame,
                merkmale: list[str] | None = None) -> np.ndarray:
    import statsmodels.api as sm

    spalten = MERKMALE if merkmale is None else list(merkmale)
    X_tr = sm.add_constant(train[spalten].astype(float), has_constant="add")
    X_te = sm.add_constant(test[spalten].astype(float),  has_constant="add")
    y_tr = train[ZIELGROESSE].astype(float)
    off_tr = np.log(train[EXPOSURE_ROH].astype(float))
    off_te = np.log(test[EXPOSURE_ROH].astype(float))

    modell = sm.GLM(y_tr, X_tr, family=sm.families.Poisson(),
                    offset=off_tr).fit()
    return np.asarray(modell.predict(X_te, offset=off_te))
\end{lstlisting}


% -------------------------------------------------------------------------
% L8  OPTIONAL, nur falls das Seitenbudget es hergibt.
%     Zeigt, dass die 19 Pruefungen keine Dekoration sind.
%     ACHTUNG: Diese Datei enthaelt echte Umlaute -> literate-Zuordnung im
%     Preamble ist hier zwingend, sonst verschwinden sie im Satz.
% -------------------------------------------------------------------------
\begin{lstlisting}[style=pythonstil,
  label={lst:tests},
  caption={Zwei der neunzehn Prüfungen an den erzeugten Dateien
           (\texttt{tests/test\_aufbereitung.py}, Z.~255--260 und 277--285,
           gekürzt)}]
def test_lags_nicht_gegenwartsbezogen():
    """Gegenprobe gegen ein vergessenes shift(): lag_1 darf nicht der Istwert sein."""
    d = regression()
    anteil = (d["lag_1"] == d["anzahl_einsaetze"]).mean()
    assert anteil < 0.10, \
        f"lag_1 stimmt in {anteil:.1%} der Fälle mit dem Istwert überein"


def test_keine_ergebnisvariablen():
    gefunden = [c for c in ERGEBNISVARIABLEN if c in klassifikation().columns]
    assert not gefunden, f"Ergebnisvariable(n) im Datensatz: {gefunden}"
\end{lstlisting}
